%An example is shown at \figref{dataflow-feature-body}.

\begin{figure}
    \begin{tikzpicture}[node distance=20pt, every node/.style={fill=white, font=\sffamily, scale=0.8}, align=left]
        \node (start) [code] {
            \begin{lstlisting}
class
    EXAMPLE
feature
    min (x, y: INTEGER): INTEGER
        -- Header comment
        do
            Result := x
        ensure
            res: (Result = x or Result = y)
                and Result <= x and Result <= y
        end
end
\end{lstlisting}
        };
        \node (start-description) [description, right of=start, xshift=130pt] {Initial feature in class};
        \draw[->] (start-description) -- (start);

        \node (prompt) [code, below of=start, yshift=-100pt] {
            \begin{lstlisting}
-- <task>
-- <model current class>
-- <model argument classes>
-- <precondition available ids>
-- <postcondition available ids>
min (x, y: INTEGER): INTEGER
    -- Header comment
    do
        <ADD BODY>
    ensure
        res: (Result = x or Result = y)
            and Result <= x and Result <= y
    end
\end{lstlisting}
        };
        \draw[->, thick] (start) -- (prompt);
        \node (prompt-description) [description, right of=prompt, xshift=130pt] {Generated prompt};
        \draw[->] (prompt-description) -- (prompt);

        \node (llm-candidate) [code, below of=prompt, yshift=-100pt] {
            \begin{lstlisting}
min (x, y: INTEGER): INTEGER
    -- Header comment
    do
        if x < y then
            Result := x
        else
            Result := y
        end
    ensure
        res: (Result = x or Result = y)
            and Result <= x and Result <= y
    end
\end{lstlisting}
        };
        \draw[->, thick] (prompt) -- (llm-candidate);
        \node (llm-candidate-description) [description, right of=llm-candidate, xshift=140pt] {Generated llm-candidate};
        \draw[->] (llm-candidate-description) -- (llm-candidate);

        \node (extracted-body) [code, below of=llm-candidate, yshift=-70pt] {
            \begin{lstlisting}
if x < y then
    Result := x
else
    Result := y
end
\end{lstlisting}
        };
        \draw[->, thick] (llm-candidate) -- (extracted-body);
        \node (extracted-body-description) [description, right of=extracted-body, xshift=120pt] {Extracted body};
        \draw[->] (extracted-body-description) -- (extracted-body);

        \node (mid-artifact) [code, below of=extracted-body, yshift=-90pt] {
            \begin{lstlisting}
class
    LLM_INSTRUMENTED_EXAMPLE
inherit
    EXAMPLE redefine min end
feature
    min (x, y: INTEGER): INTEGER
        do
            if x < y then
                Result := x
            else
                Result := y
            end
        ensure
            res: (Result = x or Result = y)
                and Result <= x and Result <= y
        end
end
\end{lstlisting}
        };
        \draw[->, thick] (extracted-body) -- (mid-artifact);
        \node (mid-artifact-description) [description, right of=mid-artifact, xshift=160pt] {Instrumented feature redefinition};
        \draw[->] (mid-artifact-description) -- (mid-artifact);

        \node (verification-result) [code, below of=mid-artifact, yshift=-80pt] { Verification output};
        \draw[->, thick] (mid-artifact) -- (verification-result);
        \draw[->, thick] (verification-result.west) -| ([xshift=-30pt] prompt.west)
            node[near start, above] {Fails} -- (prompt.west);
        \node (initial-feature-after-edits) [code, below of=verification-result, yshift=-70pt] {
            \begin{lstlisting}
min (x, y: INTEGER): INTEGER
    -- Header comment
    require
        True
    do
        if x < y then
            Result := x
        else
            Result := y
        end
    ensure
        True
        res: (Result = x or Result = y)
            and Result <= x and Result <= y
    end
\end{lstlisting}
        };
        \draw[->, thick] (verification-result) -- (initial-feature-after-edits)
            node[midway, right] {Succeds};
        \node (initial-feature-after-edits-description) [description, right of=initial-feature-after-edits, xshift=130pt] {Feature after edits};
        \draw[->] (initial-feature-after-edits-description) -- (initial-feature-after-edits);
    \end{tikzpicture}
    \caption{Dataflow of a feature in the fix feature body command}\label{dataflow-feature-body}
\end{figure}